\chapter{Projet de synth\`ese}
\label{ch:projet-final}

\begin{quote}
\emph{<<\,La meilleure fa\c{c}on d'apprendre l'analyse de donn\'ees, c'est d'analyser des donn\'ees. Pas des donn\'ees jouets --- des donn\'ees r\'eelles, avec un vrai d\'esordre, sur de vraies questions de sant\'e.\,>>}
\end{quote}

% ============================================================
\section{Consignes du projet}

Le projet de synth\`ese est votre opportunit\'e d'appliquer toutes les comp\'etences de ce cours \`a une vraie question de sant\'e. Vous travaillerez individuellement ou en bin\^omes.

\subsection{Calendrier}

\begin{itemize}
  \item \textbf{Semaine 1.} Choisir un projet, t\'el\'echarger les donn\'ees, effectuer l'exploration initiale.
  \item \textbf{Semaine 2.} Compl\'eter l'analyse, g\'en\'erer les figures et tableaux, r\'ediger le rapport.
  \item \textbf{Semaine 3.} Finaliser le rapport et pr\'eparer une pr\'esentation orale de 5 minutes.
\end{itemize}

\subsection{Attendus}

Votre projet doit inclure :
\begin{enumerate}
  \item \textbf{Une question de sant\'e claire.} Pas <<\,explorer les donn\'ees\,>> mais <<\,la pr\'evalence du diab\`ete est-elle associ\'ee \`a l'IMC apr\`es ajustement sur l'\'age et le sexe ?\,>>
  \item \textbf{Des donn\'ees publiques r\'eelles.} Pas de jeux de donn\'ees simul\'es. Toutes les sources de donn\'ees doivent \^etre cit\'ees.
  \item \textbf{Du code reproductible.} Un notebook Jupyter qui s'ex\'ecute du d\'ebut \`a la fin sans erreur.
  \item \textbf{Au moins 3 visualisations.} Pertinentes, \'etiquet\'ees et interpr\'etables.
  \item \textbf{Au moins 1 test statistique ou mod\`ele.} Avec une interpr\'etation ad\'equate.
  \item \textbf{Un rapport \'ecrit.} 5 \`a 8 pages suivant le mod\`ele ci-dessous.
  \item \textbf{Une pr\'esentation orale.} 5 minutes avec diapositives.
\end{enumerate}

\begin{warningbox}
Int\'egrit\'e acad\'emique : vous pouvez utiliser des outils d'IA (ChatGPT, GitHub Copilot) pour aider \`a \'ecrire du code, mais vous devez comprendre chaque ligne. Lors de la pr\'esentation orale, on vous demandera d'expliquer votre code et d'interpr\'eter vos r\'esultats. Copier-coller du code que vous ne comprenez pas sera \'evident.
\end{warningbox}

% ============================================================
\section{Projets sugg\'er\'es}

Choisissez l'un des cinq projets ci-dessous, ou proposez le v\^otre (sous r\'eserve de l'approbation de l'enseignant). Chaque projet sp\'ecifie l'objectif, le(s) jeu(x) de donn\'ees, les analyses requises et les livrables attendus.

% ------------------------------------------------------------
\subsection{Projet 1 : Analyse des facteurs de risque du diab\`ete}

\textbf{Objectif.} Identifier les facteurs de risque les plus importants du diab\`ete de type 2 et construire un mod\`ele pr\'edictif.

\textbf{Jeux de donn\'ees.}
\begin{itemize}
  \item \textbf{Pima Indians Diabetes Dataset} (UCI/Kaggle) :\\
  \url{https://raw.githubusercontent.com/jbrownlee/Datasets/master/pima-indians-diabetes.data.csv}\\
  768 patientes, 8 variables (grossesses, glyc\'emie, pression art\'erielle, \'epaisseur cutan\'ee, insuline, IMC, fonction de p\'edigree du diab\`ete, \'age), issue binaire.
  \item \textbf{CDC BRFSS} (Behavioral Risk Factor Surveillance System) :\\
  \url{https://www.cdc.gov/brfss/annual_data/annual_data.htm}\\
  Enqu\^ete annuelle de plus de 400\,000 adultes am\'ericains. T\'el\'echargez le CSV de l'ann\'ee la plus r\'ecente.
\end{itemize}

\textbf{Analyses requises.}
\begin{enumerate}
  \item Charger et nettoyer le jeu de donn\'ees Pima. G\'erer les z\'eros dans glyc\'emie, pression art\'erielle, IMC (ce sont des valeurs manquantes cod\'ees en 0).
  \item Analyse exploratoire : distributions, corr\'elations, bo\^ites \`a moustaches par statut diab\'etique.
  \item R\'egression logistique : quelles variables sont des pr\'edicteurs significatifs ? Rapporter les rapports des cotes avec IC \`a 95\,\%.
  \item Comparer la r\'egression logistique avec un classifieur par for\^et al\'eatoire. Rapporter exactitude, sensibilit\'e, sp\'ecificit\'e et AUC-ROC.
  \item (Bonus) R\'ep\'eter l'analyse des facteurs de risque sur les donn\'ees BRFSS et comparer les r\'esultats.
\end{enumerate}

\textbf{Livrables.}
\begin{itemize}
  \item Carte thermique de corr\'elation de toutes les variables.
  \item Courbes ROC comparant r\'egression logistique et for\^et al\'eatoire.
  \item Tableau des rapports des cotes avec intervalles de confiance.
  \item Interpr\'etation \'ecrite : quels facteurs de risque sont les plus importants cliniquement ?
\end{itemize}

\begin{clinicalnote}
Le jeu de donn\'ees Pima Indians est historiquement important mais soul\`eve des consid\'erations \'ethiques. Les donn\'ees ont \'et\'e collect\'ees aupr\`es de la communaut\'e Akimel O'odham (Pima), qui pr\'esente l'un des taux les plus \'elev\'es de diab\`ete de type 2 au monde. Lors de la pr\'esentation, reconnaissez la population et \'evitez de g\'en\'eraliser les r\'esultats \`a d'autres groupes.
\end{clinicalnote}

% ------------------------------------------------------------
\subsection{Projet 2 : Analyse de l'impact de la COVID-19 dans les pays africains}

\textbf{Objectif.} Analyser comment la COVID-19 a affect\'e diff\'eremment les pays africains et explorer les corr\'elations avec la capacit\'e du syst\`eme de sant\'e et les facteurs socio-\'economiques.

\textbf{Jeux de donn\'ees.}
\begin{itemize}
  \item \textbf{Donn\'ees JHU CSSE COVID-19} :\\
  \url{https://github.com/CSSEGISandData/COVID-19/tree/master/csse_covid_19_data/csse_covid_19_time_series}\\
  Cas confirm\'es, d\'ec\`es et gu\'erisons mondiaux (s\'eries temporelles quotidiennes).
  \item \textbf{Our World in Data} (vaccination, d\'epistage) :\\
  \url{https://raw.githubusercontent.com/owid/covid-19-data/master/public/data/owid-covid-data.csv}\\
  Jeu de donn\'ees complet avec cas, d\'ec\`es, vaccinations, d\'epistage, donn\'ees hospitali\`eres et indicateurs d\'emographiques.
  \item \textbf{Banque mondiale} (PIB, d\'epenses de sant\'e, lits d'h\^opital) :\\
  \url{https://data.worldbank.org/}\\
  Codes d'indicateurs : \texttt{SH.XPD.CHEX.PC.CD} (d\'epenses de sant\'e par habitant), \texttt{SH.MED.BEDS.ZS} (lits d'h\^opital pour 1\,000), \texttt{NY.GDP.PCAP.CD} (PIB par habitant).
  \item \textbf{Africa CDC} :\\
  \url{https://africacdc.org/covid-19/}\\
  Tableaux de bord et rapports au niveau continental.
\end{itemize}

\textbf{Analyses requises.}
\begin{enumerate}
  \item S\'electionner 10 \`a 15 pays africains de diff\'erentes sous-r\'egions (Afrique de l'Ouest, de l'Est, australe, du Nord, centrale).
  \item Tracer les courbes \'epid\'emiques (moyenne glissante 7 jours) pour chaque pays. Identifier les vagues.
  \item Calculer les indicateurs cl\'es : total des cas par million, d\'ec\`es par million, taux de l\'etalit\'e, taux de vaccination.
  \item Fusionner avec les indicateurs de la Banque mondiale. Calculer les corr\'elations entre les r\'esultats COVID-19 et la capacit\'e du syst\`eme de sant\'e (d\'epenses de sant\'e, lits d'h\^opital, densit\'e de m\'edecins).
  \item Construire une carte choropl\`ethe des taux de l\'etalit\'e \`a travers l'Afrique.
  \item (Bonus) Ex\'ecuter une r\'egression multiple : les d\'epenses de sant\'e par habitant pr\'edisent-elles le taux de l\'etalit\'e apr\`es ajustement sur l'\'age m\'edian et le PIB ?
\end{enumerate}

\textbf{Livrables.}
\begin{itemize}
  \item Courbes \'epid\'emiques pour tous les pays s\'electionn\'es (petits multiples ou graphique group\'e).
  \item Nuage de points des d\'epenses de sant\'e vs taux de l\'etalit\'e.
  \item Carte choropl\`ethe de l'Afrique montrant la charge COVID-19.
  \item Tableau r\'ecapitulatif avec les indicateurs cl\'es pour tous les pays.
\end{itemize}

% ------------------------------------------------------------
\subsection{Projet 3 : D\'eterminants de la mortalit\'e maternelle}

\textbf{Objectif.} \'Etudier quels facteurs du syst\`eme de sant\'e et socio-\'economiques sont associ\'es \`a la mortalit\'e maternelle \`a travers les pays.

\textbf{Jeux de donn\'ees.}
\begin{itemize}
  \item \textbf{GHO de l'OMS} --- Ratio de mortalit\'e maternelle (RMM) :\\
  \url{https://apps.who.int/gho/athena/api/GHO/MDG_0000000026?format=csv}
  \item \textbf{GHO de l'OMS} --- Accouchements assist\'es par du personnel qualifi\'e :\\
  \url{https://apps.who.int/gho/athena/api/GHO/MDG_0000000025?format=csv}
  \item \textbf{GHO de l'OMS} --- Couverture en soins pr\'enatals (4+ visites) :\\
  \url{https://apps.who.int/gho/athena/api/GHO/WHS4_154?format=csv}
  \item \textbf{Banque mondiale} --- Taux d'alphab\'etisation des femmes, PIB par habitant, d\'epenses de sant\'e :\\
  \url{https://data.worldbank.org/}
  \item \textbf{Programme DHS} (Enqu\^etes d\'emographiques et de sant\'e) --- indicateurs au niveau des enqu\^etes :\\
  \url{https://www.statcompiler.com/}\\
  Fournit des donn\'ees infranationales sur la sant\'e maternelle, l'utilisation de la contraception et les soins pr\'enatals pour les pays d'Afrique et d'Asie.
\end{itemize}

\textbf{Analyses requises.}
\begin{enumerate}
  \item Charger les donn\'ees de RMM et filtrer l'ann\'ee la plus r\'ecente. Identifier les 20 pays avec le RMM le plus \'elev\'e.
  \item Fusionner avec les donn\'ees d'accouchements assist\'es, de couverture en soins pr\'enatals et d'indicateurs socio-\'economiques.
  \item Analyse exploratoire : nuages de points du RMM vs chaque pr\'edicteur. Calculer les corr\'elations de Pearson et de Spearman.
  \item R\'egression lin\'eaire multiple : quels facteurs pr\'edisent significativement le RMM ? V\'erifier les hypoth\`eses du mod\`ele (graphiques des r\'esidus, normalit\'e).
  \item Cr\'eer une carte choropl\`ethe du RMM \`a travers l'Afrique.
  \item (Bonus) Comparer les tendances du RMM au fil du temps (2000--2020) pour 5 pays ayant mis en place des programmes sp\'ecifiques de sant\'e maternelle.
\end{enumerate}

\textbf{Livrables.}
\begin{itemize}
  \item Nuages de points avec droites de r\'egression (RMM vs accouchements assist\'es, RMM vs d\'epenses de sant\'e).
  \item Tableau de r\'egression avec coefficients, erreurs standard, valeurs $p$ et $R^2$.
  \item Carte choropl\`ethe de la mortalit\'e maternelle en Afrique.
  \item Discussion des implications politiques : quelles interventions pourraient r\'eduire le RMM ?
\end{itemize}

\begin{clinicalnote}
La mortalit\'e maternelle est l'une des in\'egalit\'es de sant\'e les plus criantes au monde. L'Afrique subsaharienne repr\'esente environ 70\,\% des d\'ec\`es maternels mondiaux. La plupart de ces d\'ec\`es sont \'evitables gr\^ace \`a du personnel qualifi\'e \`a l'accouchement, des soins obst\'etricaux d'urgence et un nombre ad\'equat de visites pr\'enatales.
\end{clinicalnote}

% ------------------------------------------------------------
\subsection{Projet 4 : Comparaison de mod\`eles de pr\'ediction des maladies cardiaques}

\textbf{Objectif.} Comparer plusieurs mod\`eles d'apprentissage automatique pour la pr\'ediction des maladies cardiaques et identifier les pr\'edicteurs cliniques les plus importants.

\textbf{Jeu de donn\'ees.}
\begin{itemize}
  \item \textbf{UCI Heart Disease Dataset} (Cleveland) :\\
  \url{https://archive.ics.uci.edu/ml/machine-learning-databases/heart-disease/processed.cleveland.data}\\
  303 patients, 13 variables cliniques (\'age, sexe, type de douleur thoracique, PA au repos, cholest\'erol, glyc\'emie \`a jeun, ECG au repos, fr\'equence cardiaque max, angor d'effort, d\'epression ST, pente, nombre de vaisseaux, thalass\'emie), cible binaire (pr\'esence de maladie cardiaque).
  \item Noms des colonnes : \texttt{age, sex, cp, trestbps, chol, fbs, restecg, thalach, exang, oldpeak, slope, ca, thal, target}.
\end{itemize}

\textbf{Analyses requises.}
\begin{enumerate}
  \item Charger et nettoyer les donn\'ees. G\'erer les valeurs manquantes (cod\'ees comme \texttt{?} dans le fichier original).
  \item Analyse exploratoire : distributions de chaque variable selon le statut de maladie cardiaque.
  \item S\'eparation entra\^inement-test (80/20). Utiliser une s\'eparation stratifi\'ee pour maintenir l'\'equilibre des classes.
  \item Entra\^iner et \'evaluer 4 mod\`eles :
  \begin{itemize}
    \item R\'egression logistique
    \item Arbre de d\'ecision
    \item For\^et al\'eatoire
    \item K plus proches voisins (KNN)
  \end{itemize}
  \item Pour chaque mod\`ele : rapporter exactitude, pr\'ecision, rappel, score F1 et AUC-ROC.
  \item Tracer les courbes ROC des 4 mod\`eles sur la m\^eme figure.
  \item Importance des variables : quelles variables cliniques sont les pr\'edicteurs les plus puissants ?
\end{enumerate}

\textbf{Livrables.}
\begin{itemize}
  \item Tableau comparatif des m\'etriques de performance des mod\`eles.
  \item Courbes ROC de tous les mod\`eles sur un seul graphique.
  \item Diagramme en barres de l'importance des variables (de la for\^et al\'eatoire).
  \item Matrice de confusion du mod\`ele le plus performant.
  \item Interpr\'etation clinique : les principaux pr\'edicteurs sont-ils en accord avec les recommandations en cardiologie ?
\end{itemize}

\begin{pythontip}
Utilisez \texttt{sklearn.model\_selection.cross\_val\_score} avec une validation crois\'ee \`a 5 plis au lieu d'une seule s\'eparation entra\^inement-test. Cela donne une estimation plus fiable de la performance du mod\`ele, surtout avec un petit jeu de donn\'ees comme celui-ci (303 patients).
\end{pythontip}

% ------------------------------------------------------------
\subsection{Projet 5 : Charge du paludisme et efficacit\'e des interventions}

\textbf{Objectif.} Analyser les tendances de la charge du paludisme dans les pays africains et \'evaluer si les interventions cl\'es (moustiquaires, pulv\'erisation intradomiciliaire, traitement) sont associ\'ees \`a une baisse de l'incidence.

\textbf{Jeux de donn\'ees.}
\begin{itemize}
  \item \textbf{GHO de l'OMS} --- Incidence du paludisme (pour 1\,000 personnes \`a risque) :\\
  \url{https://apps.who.int/gho/athena/api/GHO/MALARIA_INCIDENCE?format=csv}
  \item \textbf{GHO de l'OMS} --- Taux de mortalit\'e du paludisme :\\
  \url{https://apps.who.int/gho/athena/api/GHO/MALARIA_DEATHS_PER100000?format=csv}
  \item \textbf{Rapport mondial sur le paludisme de l'OMS} --- couverture des interventions (utilisation des MII, couverture PIR, traitement ACT) :\\
  \url{https://www.who.int/teams/global-malaria-programme/reports/world-malaria-report}\\
  Les tableaux de donn\'ees annexes sont disponibles en t\'el\'echargement Excel.
  \item \textbf{The Malaria Atlas Project} :\\
  \url{https://malariaatlas.org/}\\
  Estimations de pr\'evalence infranationales et cartes raster.
  \item \textbf{Programme DHS} :\\
  \url{https://www.statcompiler.com/}\\
  Donn\'ees d'enqu\^etes sur la possession et l'utilisation de MII, la couverture PIR et le recours aux soins.
\end{itemize}

\textbf{Analyses requises.}
\begin{enumerate}
  \item Charger les donn\'ees d'incidence et de mortalit\'e du paludisme. Filtrer les pays d'Afrique subsaharienne.
  \item Tracer les tendances d'incidence (2000--2022) pour les 10 pays les plus touch\'es.
  \item Fusionner avec les donn\'ees de couverture des interventions (utilisation de MII, PIR, acc\`es aux ACT).
  \item Calculer les corr\'elations entre les changements de couverture des interventions et les changements d'incidence au fil du temps (c'est-\`a-dire : l'augmentation de l'utilisation des MII pr\'edit-elle une diminution de l'incidence ?).
  \item Cr\'eer une carte choropl\`ethe de l'incidence actuelle du paludisme en Afrique.
  \item R\'egression lin\'eaire : pr\'edire le changement d'incidence du paludisme \`a partir du changement de couverture des interventions, en ajustant sur le PIB par habitant et l'urbanisation.
  \item (Bonus) Cr\'eer une comparaison avant/apr\`es pour les pays ayant intensifi\'e la distribution de MII apr\`es les subventions du Fonds mondial.
\end{enumerate}

\textbf{Livrables.}
\begin{itemize}
  \item Graphique de s\'eries temporelles de l'incidence du paludisme pour 10 pays (2000--2022).
  \item Nuage de points du changement d'utilisation de MII vs changement d'incidence avec droite de r\'egression.
  \item Carte choropl\`ethe de l'incidence du paludisme en Afrique.
  \item Tableau de r\'egression avec interpr\'etation.
  \item Discussion politique : quelles interventions montrent l'association la plus forte avec la baisse du paludisme ?
\end{itemize}

\begin{clinicalnote}
Entre 2000 et 2015, la mortalit\'e due au paludisme en Afrique a diminu\'e d'environ 44\,\%, largement gr\^ace \`a l'intensification des MII, de la PIR et des th\'erapies combin\'ees \`a base d'art\'emisinine (ACT). Depuis 2015, les progr\`es stagnent, ce qui rend cette analyse opportune et pertinente pour les politiques de sant\'e.
\end{clinicalnote}

% ============================================================
\section{Mod\`ele de rapport}

Votre rapport \'ecrit doit suivre cette structure (5 \`a 8 pages, hors annexe de code) :

\subsection{1. Introduction (0,5 \`a 1 page)}

\begin{itemize}
  \item Contexte sanitaire : pourquoi cette question est-elle importante ?
  \item Question de recherche ou objectif sp\'ecifique.
  \item Br\`eve description du (des) jeu(x) de donn\'ees utilis\'e(s).
\end{itemize}

\subsection{2. Description des donn\'ees (1 page)}

\begin{itemize}
  \item Source(s) et citation(s).
  \item Nombre d'observations et de variables.
  \item Variables cl\'es : nom, type, unit\'e et signification clinique.
  \item Tableau de statistiques descriptives (\texttt{.describe()}).
  \item Donn\'ees manquantes : combien, quelles variables, comment trait\'ees.
\end{itemize}

\subsection{3. M\'ethodes (1 page)}

\begin{itemize}
  \item \'Etapes de nettoyage (gestion des valeurs manquantes, recodage des variables, traitement des valeurs aberrantes).
  \item Tests statistiques utilis\'es et pourquoi (par ex.\ <<\,nous avons utilis\'e le test du chi-deux car les deux variables sont cat\'egorielles\,>>).
  \item Mod\`eles utilis\'es (r\'egression logistique, for\^et al\'eatoire, etc.) avec hyperparam\`etres.
  \item Strat\'egie de validation (s\'eparation entra\^inement-test, validation crois\'ee).
\end{itemize}

\subsection{4. R\'esultats (1,5 \`a 2 pages)}

\begin{itemize}
  \item Figures et tableaux avec l\'egendes.
  \item Rapporter les statistiques de test, valeurs $p$, intervalles de confiance.
  \item M\'etriques de performance du mod\`ele.
  \item Ne pas interpr\'eter ici --- pr\'esentez seulement les chiffres.
\end{itemize}

\subsection{5. Discussion (1 \`a 1,5 page)}

\begin{itemize}
  \item Que signifient les r\'esultats cliniquement ?
  \item Comment vos r\'esultats se comparent-ils \`a la litt\'erature publi\'ee ?
  \item Qu'est-ce qui vous a surpris ?
  \item Quelles sont les implications pratiques ?
\end{itemize}

\subsection{6. Limites (0,5 page)}

\begin{itemize}
  \item Probl\`emes de qualit\'e des donn\'ees (donn\'ees manquantes, erreurs de mesure, biais de s\'election).
  \item G\'en\'eralisabilit\'e : le jeu de donn\'ees repr\'esente-t-il la population qui vous int\'eresse ?
  \item Quelles analyses feriez-vous avec plus de temps ou de donn\'ees ?
\end{itemize}

\begin{warningbox}
Chaque figure doit avoir un titre, des axes \'etiquet\'es et une l\'egende expliquant ce que le lecteur doit y voir. Chaque tableau doit avoir une ligne d'en-t\^ete et une l\'egende. Les figures sans \'etiquettes sont la raison la plus fr\'equente de perte de points.
\end{warningbox}

% ============================================================
\section{Consignes de pr\'esentation}

Vous ferez une pr\'esentation orale de 5 minutes suivie de 2 \`a 3 minutes de questions.

\subsection{Structure}

\begin{enumerate}
  \item \textbf{Diapositive 1 : Titre.} Titre du projet, votre nom, date.
  \item \textbf{Diapositive 2 : Motivation.} Pourquoi cette question de sant\'e est-elle importante ? Un ou deux faits cl\'es.
  \item \textbf{Diapositive 3 : Donn\'ees.} Quelles donn\'ees avez-vous utilis\'ees ? Combien d'observations ? Quelles sont les variables cl\'es ?
  \item \textbf{Diapositives 4--5 : R\'esultats cl\'es.} Vos 2 ou 3 figures ou tableaux les plus importants. Guidez le public \`a travers chacun.
  \item \textbf{Diapositive 6 : Conclusions.} Un ou deux messages \`a retenir. Que doit retenir le public ?
\end{enumerate}

\subsection{Conseils}

\begin{itemize}
  \item \textbf{5 minutes, c'est court.} Entra\^inez-vous avec un chronom\`etre. \'Eliminez tout ce qui n'est pas essentiel.
  \item \textbf{Un message par diapositive.} Ne surchargez pas les diapositives de texte.
  \item \textbf{Parlez au public, pas \`a l'\'ecran.}
  \item \textbf{Anticipez les questions.} Connaissez vos limites et soyez pr\^et \`a en discuter.
  \item \textbf{Pas de code sur les diapositives.} Montrez les r\'esultats, pas le code qui les a produits.
\end{itemize}

\begin{pythontip}
Sauvegardez vos figures en PNG haute r\'esolution pour les diapositives :
\begin{minted}{python}
fig.savefig("figure1.png", dpi=300, bbox_inches="tight")
\end{minted}
Cela \'evite les captures d'\'ecran floues et garantit que vos graphiques sont professionnels dans une pr\'esentation.
\end{pythontip}

% ============================================================
\section{Bar\`eme de notation}

Le projet de synth\`ese est not\'e sur 100 points :

\begin{center}
\begin{tabular}{lcp{8cm}}
\toprule
\textbf{Composante} & \textbf{Points} & \textbf{Crit\`eres} \\
\midrule
Question de recherche & 10 & Claire, sp\'ecifique et pertinente pour la sant\'e \\
Gestion des donn\'ees & 15 & Chargement, nettoyage, traitement des valeurs manquantes et documentation ad\'equats \\
Analyse exploratoire & 15 & Statistiques descriptives, distributions et au moins 3 visualisations bien \'etiquet\'ees \\
Analyse statistique & 20 & Choix et application corrects des tests ou mod\`eles ; interpr\'etation ad\'equate des valeurs $p$, intervalles de confiance et tailles d'effet \\
Qualit\'e du rapport & 20 & Suit le mod\`ele ; r\'edaction claire ; figures avec titres, \'etiquettes et l\'egendes ; limites honn\^etement discut\'ees \\
Qualit\'e du code & 10 & Notebook reproductible qui s'ex\'ecute du d\'ebut \`a la fin ; commentaires expliquant les \'etapes cl\'es ; pas de code inutile \\
Pr\'esentation orale & 10 & Claire, dans la limite de temps, r\'epond aux questions avec assurance \\
\midrule
\textbf{Total} & \textbf{100} & \\
\bottomrule
\end{tabular}
\end{center}

\subsection{Limites de notes}

\begin{itemize}
  \item \textbf{A (90--100) :} Excellente analyse avec interpr\'etation perspicace. Code propre et bien document\'e. Va au-del\`a des exigences minimales.
  \item \textbf{B (75--89) :} Analyse solide avec m\'ethodologie correcte. Probl\`emes mineurs de pr\'esentation ou d'interpr\'etation.
  \item \textbf{C (60--74) :} R\'epond aux exigences de base mais pr\'esente des lacunes notables (par ex.\ visualisations manquantes, interpr\'etation superficielle, erreurs de code).
  \item \textbf{D/F ($<60$) :} Analyse incompl\`ete, erreurs m\'ethodologiques majeures ou code non reproductible.
\end{itemize}

% ============================================================
\section{Pour d\'emarrer : une liste de v\'erification}

\begin{exercise}
Avant de partir aujourd'hui, compl\'etez cette liste de v\'erification :
\begin{enumerate}
  \item Choisissez votre projet (1 \`a 5 ou une proposition personnalis\'ee).
  \item T\'el\'echargez votre jeu de donn\'ees principal. Chargez-le dans un notebook Jupyter et ex\'ecutez \texttt{.head()}, \texttt{.shape}, \texttt{.info()} et \texttt{.describe()}.
  \item R\'edigez votre question de recherche en une phrase.
  \item Listez 3 visualisations que vous pr\'evoyez de cr\'eer.
  \item Listez 1 test statistique ou mod\`ele que vous pr\'evoyez d'utiliser.
  \item Identifiez les probl\`emes potentiels : donn\'ees manquantes, noms de colonnes d\'esordonn\'es, multiples fichiers \`a fusionner.
  \item Cr\'eez un dossier de projet avec la structure :
\begin{minted}{text}
projet_synthese/
    data/
        raw/          # fichiers originaux telecharges
        cleaned/      # donnees traitees
    notebooks/
        01_exploration.ipynb
        02_analyse.ipynb
    figures/
    rapport.pdf
\end{minted}
  \item Soumettez votre choix de projet et votre question de recherche \`a l'enseignant avant la fin de la semaine.
\end{enumerate}
\end{exercise}

% ============================================================
\section{R\'esum\'e du chapitre}

\begin{itemize}
  \item Le projet de synth\`ese int\`egre toutes les comp\'etences du cours : chargement, nettoyage, exploration, visualisation, analyse statistique et interpr\'etation des donn\'ees.
  \item Cinq projets sugg\'er\'es couvrent le diab\`ete, la COVID-19, la mortalit\'e maternelle, les maladies cardiaques et le paludisme --- tous utilisant des jeux de donn\'ees publics et r\'eels.
  \item Chaque projet exige une question de sant\'e claire, du code reproductible, au moins 3 visualisations et au moins 1 test statistique ou mod\`ele.
  \item Le rapport suit une structure standard : Introduction, Donn\'ees, M\'ethodes, R\'esultats, Discussion, Limites.
  \item Les pr\'esentations durent 5 minutes --- concentrez-vous sur la motivation, les r\'esultats cl\'es et les conclusions.
  \item Le bar\`eme de notation r\'ecompense la clart\'e de la r\'eflexion et l'honn\^etet\'e de l'interpr\'etation plut\^ot que la complexit\'e m\'ethodologique.
\end{itemize}
